\documentclass[pdflatex,sn-mathphys-num]{sn-jnl}

\usepackage{graphicx}%
\usepackage{multirow}%
\usepackage{amsmath,amssymb,amsfonts}%
\usepackage{amsthm}%
\usepackage{mathrsfs}%
\usepackage[title]{appendix}%
\usepackage{xcolor}%
\usepackage{textcomp}%
\usepackage{manyfoot}%
\usepackage{booktabs}%
\usepackage{algorithm}%
\usepackage{algorithmicx}%
\usepackage{algpseudocode}%
\usepackage{listings}%
\usepackage{lmodern}%
\usepackage{siunitx}
\usepackage{microtype}
\usepackage{derivative}
\usepackage[version=4]{mhchem}
\usepackage{cleveref}
\usepackage[spanish]{babel}

\renewcommand{\arraystretch}{1.2}
\sisetup{group-digits=true,
  group-separator={\,},
  separate-uncertainty
}

\theoremstyle{thmstyleone}%
\newtheorem{theorem}{Teorema}%
\newtheorem{proposition}[theorem]{Proposición}%

\theoremstyle{thmstyletwo}%
\newtheorem{example}{Ejemplo}%
\newtheorem{remark}{Observación}%

\theoremstyle{thmstylethree}%
\newtheorem{definition}{Definición}%

\DeclareSIUnit\angstrom{\text{\r{A}}}

\raggedbottom
%%\unnumbered% uncomment this for unnumbered level heads

\begin{document}

\title[Resumen de Charla Nanociencia]{Nanociencia y Materiales
Avanzados: Resumen de la Charla Impartida por Dav Alessia}

\author{\fnm{Julian L.}
\sur{Avila-Martinez}}\email{jlavilam@udistrital.edu.co}

\affil{\orgdiv{Programa Académico de Física}, \orgname{Universidad Distrital
Francisco José de Caldas}, \orgaddress{\city{Bogotá D.C.},
\country{Colombia}}}

\abstract{
Resumen de la charla impartida por la investigadora Dav Alessia.
Se sintetizan los principios expuestos sobre nanomateriales, detallando
las fuerzas intermoleculares como mecanismo rector del autoensamblaje.
Se documenta su exposición sobre la aplicación de marcos
metal-orgánicos (MOF) en el suministro localizado de fármacos, las
propiedades de semiconductores a nanoescala y la estabilización
electrónica de perovskitas mediante funcionalización con derivados de
carbono.
}

\keywords{Química supramolecular, marcos metal-orgánicos, puntos cuánticos,
perovskitas, fullerenos, confinamiento cuántico}

\maketitle

\section{Fuerzas Intermoleculares y Autoensamblaje}
\label{sec:fuerzas}

Durante su intervención, Alessia estableció que la comprensión rigurosa de
la nanociencia exige distinguir entre interacciones intramoleculares e
intermoleculares. Los enlaces covalentes, iónicos y metálicos
constituyen fuerzas intramoleculares que definen la estructura intrínseca de
la molécula. Por el contrario, el diseño de arquitecturas a
nanoescala está gobernado por interacciones no covalentes, tales como
interacciones ion-dipolo, dipolo-dipolo, puentes de hidrógeno y fuerzas de
dispersión.

Estas fuerzas operan a escalas espaciales amplias; mientras un enlace
intramolecular en el \ce{H2O} exhibe una separación de
\qty{96}{\pico\meter}, las distancias intermoleculares escalan a
\qty{300}{\pico\meter}. Esta dilatación espacial impone
interacciones de menor magnitud, proveyendo la flexibilidad requerida para
las arquitecturas huésped-hospedador. El rendimiento de este
autoensamblaje se optimiza mediante la preorganización estructural y el
efecto quelato, limitando el costo energético durante la asimilación del
huésped.

\section{Marcos Metal-Orgánicos (MOF) y Terapia Dirigida}
\label{sec:mof}

La conferencista expuso el funcionamiento de los marcos metal-orgánicos como
jaulas reticulares que permiten el encapsulamiento de sustancias.
En aplicaciones biomédicas, estos cristales se emplean como
vectores oncológicos alojando un fármaco y recubriéndose con ligandos
direccionales y nanopartículas de oro. Su mecanismo explota
el gradiente de acidez tisular, catalizando la degradación del MOF en el
entorno neoplásico, el cual presenta un pH de \num{5.5} en contraste con el
pH de \num{7.4} de los tejidos sanos.

Las nanopartículas de oro, sintetizadas a partir de \ce{AuCl4-} con
diámetros ideales de \qtyrange{12}{18}{\nano\meter}, actúan
como transductores. Absorben luz visible y disipan energía
térmica en la región infrarroja, induciendo hipertermia localizada para
erradicar células comprometidas.

\section{Puntos Cuánticos y Semiconductores}
\label{sec:semiconductores}

Se describieron los nanodots como semiconductores confinados en regímenes de
\qtyrange{1}{10}{\nano\meter}.
Sus propiedades ópticas dependen de restricciones impuestas por el confinamiento
cuántico.
Al alterar dicho confinamiento, la brecha de energía sufre variaciones que
provocan desplazamientos espectrales hacia el azul en partículas pequeñas y
hacia el rojo en dimensiones mayores. Esta fenomenología permite
absorber radiación ultravioleta y reconvertirla al espectro visible para
optimizar celdas solares.

Respecto a la generación energética, Alessia detalló el uso de redes
cristalinas tipo perovskita. Para eludir la toxicidad del plomo,
la síntesis actual se orienta al estaño. La propensión del
estaño a oxidarse se compensa mediante integración con fullerenos.
Estas estructuras de carbono desplazan la densidad electrónica
hacia su superficie, estableciendo una movilidad de carga que preserva la
perovskita. Esta lógica funcional aplica igualmente al
integrar grafeno y nanotubos de carbono monocapa y multicapa.

\section{Dispositivos Moleculares}
\label{sec:dispositivos}

La presentación abordó la convergencia entre síntesis química y
mecánica fundamental para construir ensamblajes dinámicos.
Estructuras como los rotaxanos ejecutan traslaciones controladas o actúan
como interruptores ante estímulos externos como variaciones de pH, diferencia de
potencial eléctrico o luz.
De manera análoga a la transformación de ATP en la
mitocondria, la ingeniería a nanoescala ha desarrollado
vehículos rotacionales y ascensores moleculares.

\section{Conclusiones}
\label{sec:conclusion}

La conferencista concluyó que la capacidad de diseñar la materia átomo a
átomo es un pilar fundamental para el desarrollo biomédico y tecnológico
. El espectro de acción de estos materiales supramoleculares
unifica disciplinas, permitiendo reproducir funciones macroscópicas
mediante el control riguroso de la estructura a nanoescala.

\end{document}
